\IfFileExists{stacks-project.cls}{%
\documentclass{stacks-project}
}{%
\documentclass{amsart}
}

% For dealing with references we use the comment environment
\usepackage{verbatim}
\newenvironment{reference}{\comment}{\endcomment}
%\newenvironment{reference}{}{}
\newenvironment{slogan}{\comment}{\endcomment}
\newenvironment{history}{\comment}{\endcomment}

% For commutative diagrams we use Xy-pic
\usepackage[all]{xy}

% We use 2cell for 2-commutative diagrams.
\xyoption{2cell}
\UseAllTwocells

% We use multicol for the list of chapters between chapters
\usepackage{multicol}

% This is generally recommended for better output
\usepackage{lmodern}
\usepackage[T1]{fontenc}

% For cross-file-references
\usepackage{xr-hyper}

% Package for hypertext links:
\usepackage{hyperref}

% For any local file, say "hello.tex" you want to link to please
% use \externaldocument[hello-]{hello}
\externaldocument[introduction-]{introduction}
\externaldocument[conventions-]{conventions}
\externaldocument[sets-]{sets}
\externaldocument[categories-]{categories}
\externaldocument[topology-]{topology}
\externaldocument[sheaves-]{sheaves}
\externaldocument[sites-]{sites}
\externaldocument[stacks-]{stacks}
\externaldocument[fields-]{fields}
\externaldocument[algebra-]{algebra}
\externaldocument[brauer-]{brauer}
\externaldocument[homology-]{homology}
\externaldocument[derived-]{derived}
\externaldocument[simplicial-]{simplicial}
\externaldocument[more-algebra-]{more-algebra}
\externaldocument[smoothing-]{smoothing}
\externaldocument[modules-]{modules}
\externaldocument[sites-modules-]{sites-modules}
\externaldocument[injectives-]{injectives}
\externaldocument[cohomology-]{cohomology}
\externaldocument[sites-cohomology-]{sites-cohomology}
\externaldocument[dga-]{dga}
\externaldocument[dpa-]{dpa}
\externaldocument[sdga-]{sdga}
\externaldocument[hypercovering-]{hypercovering}
\externaldocument[schemes-]{schemes}
\externaldocument[constructions-]{constructions}
\externaldocument[properties-]{properties}
\externaldocument[morphisms-]{morphisms}
\externaldocument[coherent-]{coherent}
\externaldocument[divisors-]{divisors}
\externaldocument[limits-]{limits}
\externaldocument[varieties-]{varieties}
\externaldocument[topologies-]{topologies}
\externaldocument[descent-]{descent}
\externaldocument[perfect-]{perfect}
\externaldocument[more-morphisms-]{more-morphisms}
\externaldocument[flat-]{flat}
\externaldocument[groupoids-]{groupoids}
\externaldocument[more-groupoids-]{more-groupoids}
\externaldocument[etale-]{etale}
\externaldocument[chow-]{chow}
\externaldocument[intersection-]{intersection}
\externaldocument[pic-]{pic}
\externaldocument[weil-]{weil}
\externaldocument[adequate-]{adequate}
\externaldocument[dualizing-]{dualizing}
\externaldocument[duality-]{duality}
\externaldocument[discriminant-]{discriminant}
\externaldocument[derham-]{derham}
\externaldocument[local-cohomology-]{local-cohomology}
\externaldocument[algebraization-]{algebraization}
\externaldocument[curves-]{curves}
\externaldocument[resolve-]{resolve}
\externaldocument[models-]{models}
\externaldocument[functors-]{functors}
\externaldocument[equiv-]{equiv}
\externaldocument[pione-]{pione}
\externaldocument[etale-cohomology-]{etale-cohomology}
\externaldocument[proetale-]{proetale}
\externaldocument[relative-cycles-]{relative-cycles}
\externaldocument[more-etale-]{more-etale}
\externaldocument[trace-]{trace}
\externaldocument[crystalline-]{crystalline}
\externaldocument[spaces-]{spaces}
\externaldocument[spaces-properties-]{spaces-properties}
\externaldocument[spaces-morphisms-]{spaces-morphisms}
\externaldocument[decent-spaces-]{decent-spaces}
\externaldocument[spaces-cohomology-]{spaces-cohomology}
\externaldocument[spaces-limits-]{spaces-limits}
\externaldocument[spaces-divisors-]{spaces-divisors}
\externaldocument[spaces-over-fields-]{spaces-over-fields}
\externaldocument[spaces-topologies-]{spaces-topologies}
\externaldocument[spaces-descent-]{spaces-descent}
\externaldocument[spaces-perfect-]{spaces-perfect}
\externaldocument[spaces-more-morphisms-]{spaces-more-morphisms}
\externaldocument[spaces-flat-]{spaces-flat}
\externaldocument[spaces-groupoids-]{spaces-groupoids}
\externaldocument[spaces-more-groupoids-]{spaces-more-groupoids}
\externaldocument[bootstrap-]{bootstrap}
\externaldocument[spaces-pushouts-]{spaces-pushouts}
\externaldocument[spaces-chow-]{spaces-chow}
\externaldocument[groupoids-quotients-]{groupoids-quotients}
\externaldocument[spaces-more-cohomology-]{spaces-more-cohomology}
\externaldocument[spaces-simplicial-]{spaces-simplicial}
\externaldocument[spaces-duality-]{spaces-duality}
\externaldocument[formal-spaces-]{formal-spaces}
\externaldocument[restricted-]{restricted}
\externaldocument[spaces-resolve-]{spaces-resolve}
\externaldocument[formal-defos-]{formal-defos}
\externaldocument[defos-]{defos}
\externaldocument[cotangent-]{cotangent}
\externaldocument[examples-defos-]{examples-defos}
\externaldocument[algebraic-]{algebraic}
\externaldocument[examples-stacks-]{examples-stacks}
\externaldocument[stacks-sheaves-]{stacks-sheaves}
\externaldocument[criteria-]{criteria}
\externaldocument[artin-]{artin}
\externaldocument[quot-]{quot}
\externaldocument[stacks-properties-]{stacks-properties}
\externaldocument[stacks-morphisms-]{stacks-morphisms}
\externaldocument[stacks-limits-]{stacks-limits}
\externaldocument[stacks-cohomology-]{stacks-cohomology}
\externaldocument[stacks-perfect-]{stacks-perfect}
\externaldocument[stacks-introduction-]{stacks-introduction}
\externaldocument[stacks-more-morphisms-]{stacks-more-morphisms}
\externaldocument[stacks-geometry-]{stacks-geometry}
\externaldocument[moduli-]{moduli}
\externaldocument[moduli-curves-]{moduli-curves}
\externaldocument[examples-]{examples}
\externaldocument[exercises-]{exercises}
\externaldocument[guide-]{guide}
\externaldocument[desirables-]{desirables}
\externaldocument[coding-]{coding}
\externaldocument[obsolete-]{obsolete}
\externaldocument[fdl-]{fdl}
\externaldocument[index-]{index}

% Theorem environments.
%
\theoremstyle{plain}
\newtheorem{theorem}[subsection]{Theorem}
\newtheorem{proposition}[subsection]{Proposition}
\newtheorem{lemma}[subsection]{Lemma}

\theoremstyle{definition}
\newtheorem{definition}[subsection]{Definition}
\newtheorem{example}[subsection]{Example}
\newtheorem{exercise}[subsection]{Exercise}
\newtheorem{situation}[subsection]{Situation}

\theoremstyle{remark}
\newtheorem{remark}[subsection]{Remark}
\newtheorem{remarks}[subsection]{Remarks}

\numberwithin{equation}{subsection}

% Macros
%
\def\lim{\mathop{\mathrm{lim}}\nolimits}
\def\colim{\mathop{\mathrm{colim}}\nolimits}
\def\Spec{\mathop{\mathrm{Spec}}}
\def\Hom{\mathop{\mathrm{Hom}}\nolimits}
\def\Ext{\mathop{\mathrm{Ext}}\nolimits}
\def\SheafHom{\mathop{\mathcal{H}\!\mathit{om}}\nolimits}
\def\SheafExt{\mathop{\mathcal{E}\!\mathit{xt}}\nolimits}
\def\Sch{\mathit{Sch}}
\def\Mor{\mathop{\mathrm{Mor}}\nolimits}
\def\Ob{\mathop{\mathrm{Ob}}\nolimits}
\def\Sh{\mathop{\mathit{Sh}}\nolimits}
\def\NL{\mathop{N\!L}\nolimits}
\def\CH{\mathop{\mathrm{CH}}\nolimits}
\def\proetale{{pro\text{-}\acute{e}tale}}
\def\etale{{\acute{e}tale}}
\def\QCoh{\mathit{QCoh}}
\def\Ker{\mathop{\mathrm{Ker}}}
\def\Im{\mathop{\mathrm{Im}}}
\def\Coker{\mathop{\mathrm{Coker}}}
\def\Coim{\mathop{\mathrm{Coim}}}

% Boxtimes
%
\DeclareMathSymbol{\boxtimes}{\mathbin}{AMSa}{"02}

%
% Macros for moduli stacks/spaces
%
\def\QCohstack{\mathcal{QC}\!\mathit{oh}}
\def\Cohstack{\mathcal{C}\!\mathit{oh}}
\def\Spacesstack{\mathcal{S}\!\mathit{paces}}
\def\Quotfunctor{\mathrm{Quot}}
\def\Hilbfunctor{\mathrm{Hilb}}
\def\Curvesstack{\mathcal{C}\!\mathit{urves}}
\def\Polarizedstack{\mathcal{P}\!\mathit{olarized}}
\def\Complexesstack{\mathcal{C}\!\mathit{omplexes}}
% \Pic is the operator that assigns to X its picard group, usage \Pic(X)
% \Picardstack_{X/B} denotes the Picard stack of X over B
% \Picardfunctor_{X/B} denotes the Picard functor of X over B
\def\Pic{\mathop{\mathrm{Pic}}\nolimits}
\def\Picardstack{\mathcal{P}\!\mathit{ic}}
\def\Picardfunctor{\mathrm{Pic}}
\def\Deformationcategory{\mathcal{D}\!\mathit{ef}}
\usepackage{cleveref}
\crefname{originalreference}{original reference}{original references}
\Crefname{originalreference}{Original reference}{Original references}
\makeatletter
\newlabel{homology-lemma-double-complex-gives-resolution}{{0133}{source}{}{sourcepointer-1}{}}
\newlabel{homology-lemma-double-complex-gives-resolution@cref}{{[originalreference][0][]0133}{[1][0][]source}{}{}{}}
\makeatother
\crefname{originalreference}{original reference}{original references}
\Crefname{originalreference}{Original reference}{Original references}
\makeatletter
\newlabel{homology-lemma-first-quadrant-ss}{{0132}{source}{}{sourcepointer-2}{}}
\newlabel{homology-lemma-first-quadrant-ss@cref}{{[originalreference][0][]0132}{[1][0][]source}{}{}{}}
\makeatother
\crefname{originalreference}{original reference}{original references}
\Crefname{originalreference}{Original reference}{Original references}
\makeatletter
\newlabel{homology-lemma-ss-double-complex}{{0130}{source}{}{sourcepointer-3}{}}
\newlabel{homology-lemma-ss-double-complex@cref}{{[originalreference][0][]0130}{[1][0][]source}{}{}{}}
\makeatother
\title{The two spectral sequences of a Cartan-Eilenberg resolution compute R\textasciicircum{}qF(K\textasciicircum{}p) and R\textasciicircum{}pF(H\textasciicircum{}q(K))}
\author{The Stacks Project Authors}
\date{stacks-project:derived.tex:local-snapshot-20261012; GNU FDL 1.2 or later}
\begin{document}
\maketitle
\noindent\textit{Editorial note.} Counted span is the author's own proof environment, source lines 7415-7468, following the statement at 7393-7413.\par\medskip
\noindent\textit{Editorial note.} The three opening remarks are stated as 'used without further mention' but are explicitly listed with their inputs (7416-7430), recorded.\par\medskip
\noindent\textit{Allowed standard prerequisites (editorial).}\begin{itemize}
\item Total complex of a Cartan-Eilenberg resolution. homology-lemma-double-complex-gives-resolution, cross-chapter citation at proof line 7425 (tags mechanism). Supplies that Tot(I) is an injective resolution of K, so RF is defined at K with value F(Tot(I)).
\item Spectral sequences of a double complex. homology-lemma-ss-double-complex and homology-lemma-first-quadrant-ss, cross-chapter citations at proof lines 7433 and 7436 (tags mechanism). Supplies the two spectral sequences, their boundedness, convergence and finite filtrations.
\end{itemize}\medskip
\ifdefined\hypertarget\hypertarget{sourcepointer-1}{}\fi
\noindent\textit{Source reference.} Original reference 0133: focused-seeds/geometry/S9-der-two-ss-complex-functor/ref-tags.txt, source lines 5--5. Navigation pointer only: the counted proof cites this statement by \textbackslash{}ref as a standing input; the deduction the counted proof has to make is carried out inside the counted span, and no part of that deduction is delegated to this pointer.\par\medskip
\ifdefined\hypertarget\hypertarget{sourcepointer-2}{}\fi
\noindent\textit{Source reference.} Original reference 0132: focused-seeds/geometry/S9-der-two-ss-complex-functor/ref-tags.txt, source lines 6--6. Navigation pointer only: the counted proof cites this statement by \textbackslash{}ref as a standing input; the deduction the counted proof has to make is carried out inside the counted span, and no part of that deduction is delegated to this pointer.\par\medskip
\ifdefined\hypertarget\hypertarget{sourcepointer-3}{}\fi
\noindent\textit{Source reference.} Original reference 0130: focused-seeds/geometry/S9-der-two-ss-complex-functor/ref-tags.txt, source lines 7--7. Navigation pointer only: the counted proof cites this statement by \textbackslash{}ref as a standing input; the deduction the counted proof has to make is carried out inside the counted span, and no part of that deduction is delegated to this pointer.\par\medskip
\begin{lemma}
\label{lemma-two-ss-complex-functor}
Let $F : \mathcal{A} \to \mathcal{B}$ be a left exact functor of
abelian categories.
Let $K^\bullet$ be a bounded below complex of $\mathcal{A}$.
Let $I^{\bullet, \bullet}$ be a Cartan-Eilenberg resolution
for $K^\bullet$. The spectral sequences
$({}'E_r, {}'d_r)_{r \geq 0}$ and $({}''E_r, {}''d_r)_{r \geq 0}$
associated to the double complex $F(I^{\bullet, \bullet})$
satisfy the relations
$$
{}'E_1^{p, q} = R^qF(K^p)
\quad
\text{and}
\quad
{}''E_2^{p, q} = R^pF(H^q(K^\bullet))
$$
Moreover, these spectral sequences are bounded, converge to
$H^*(RF(K^\bullet))$, and the associated induced filtrations on
$H^n(RF(K^\bullet))$ are finite.
\end{lemma}

\begin{proof}
We will use the following remarks without further mention:
\begin{enumerate}
\item As $I^{p, \bullet}$ is an injective resolution of
$K^p$ we see that $RF$ is defined at $K^p[0]$
with value $F(I^{p, \bullet})$.
\item As $H^p_I(I^{\bullet, \bullet})$ is an injective resolution of
$H^p(K^\bullet)$ the derived functor $RF$ is defined at $H^p(K^\bullet)[0]$
with value $F(H^p_I(I^{\bullet, \bullet}))$.
\item By
Homology, Lemma \ref{homology-lemma-double-complex-gives-resolution}
the total complex $\text{Tot}(I^{\bullet, \bullet})$
is an injective resolution of
$K^\bullet$. Hence $RF$ is defined at $K^\bullet$ with value
$F(\text{Tot}(I^{\bullet, \bullet}))$.
\end{enumerate}
Consider the two spectral sequences associated to the double complex
$L^{\bullet, \bullet} = F(I^{\bullet, \bullet})$, see
Homology, Lemma \ref{homology-lemma-ss-double-complex}.
These are both bounded, converge to $H^*(\text{Tot}(L^{\bullet, \bullet}))$,
and induce finite filtrations on $H^n(\text{Tot}(L^{\bullet, \bullet}))$, see
Homology, Lemma \ref{homology-lemma-first-quadrant-ss}.
Since
$\text{Tot}(L^{\bullet, \bullet}) =
\text{Tot}(F(I^{\bullet, \bullet})) =
F(\text{Tot}(I^{\bullet, \bullet}))$ computes
$RF(K^\bullet)$ we find the final assertion of the lemma holds true.

\medskip\noindent
Computation of the first spectral sequence. We have
${}'E_1^{p, q} = H^q(L^{p, \bullet})$ in other words
$$
{}'E_1^{p, q} = H^q(F(I^{p, \bullet})) = R^qF(K^p)
$$
as desired. Observe for later use that the maps
${}'d_1^{p, q} : {}'E_1^{p, q} \to {}'E_1^{p + 1, q}$ are the maps
$R^qF(K^p) \to R^qF(K^{p + 1})$ induced by $K^p \to K^{p + 1}$
and the fact that $R^qF$ is a functor.

\medskip\noindent
Computation of the second spectral sequence. We have
${}''E_1^{p, q} = H^q(L^{\bullet, p}) = H^q(F(I^{\bullet, p}))$.
Note that the complex $I^{\bullet, p}$ is bounded below,
consists of injectives, and moreover each kernel, image, and
cohomology group of the differentials is an injective object
of $\mathcal{A}$. Hence we can split the differentials, i.e.,
each differential is a split surjection onto a direct summand.
It follows that the same is true after applying $F$. Hence
${}''E_1^{p, q} = F(H^q(I^{\bullet, p})) = F(H^q_I(I^{\bullet, p}))$.
The differentials on this are $(-1)^q$ times $F$ applied to
the differential of the complex $H^p_I(I^{\bullet, \bullet})$
which is an injective resolution of $H^p(K^\bullet)$. Hence the
description of the $E_2$ terms.
\end{proof}

\end{document}
